\documentclass[10.5pt]{article}
% Prototype spread: The Work That Vanishes, Chapter 1 (XeLaTeX)
\usepackage[paperwidth=7in,paperheight=10in,left=0.75in,top=0.85in,bottom=0.95in,textwidth=4.25in,marginparwidth=1.3in,marginparsep=0.15in]{geometry}
\usepackage{fontspec}
\setmainfont{DejaVu Sans}
\setsansfont{DejaVu Sans}
\setmonofont{DejaVu Sans Mono}[Scale=0.88]
\usepackage{unicode-math}
\setmathfont{DejaVu Math TeX Gyre}
\usepackage{xcolor}
\usepackage{tikz}
\usetikzlibrary{calc,decorations.pathmorphing,shapes.geometric}
\usepackage[most]{tcolorbox}
\usepackage{fancyvrb}
\usepackage{graphicx}
\usepackage{marginnote}
\usepackage{eso-pic}
\usepackage{enumitem}
\usepackage{microtype}
\usepackage{booktabs}
\setlength{\parindent}{0pt}
\setlength{\parskip}{0.45em}
\renewcommand{\footnoterule}{\kern-3pt\hrule width 2.5cm height 0.8pt\kern 3pt}
\pagestyle{empty}

% ---------- palette ----------
\definecolor{paper}{HTML}{FFF4CF}
\definecolor{inkbrown}{HTML}{7A3E12}
\definecolor{hudbg}{HTML}{06140D}
\definecolor{hudgreen}{HTML}{4DFF88}
\definecolor{hudamber}{HTML}{FFC233}
\definecolor{hudred}{HTML}{FF6A55}
\definecolor{hudorange}{HTML}{FF8A1F}
\definecolor{hudnavy}{HTML}{1F3A5F}
\definecolor{slateA}{HTML}{6F7F94}
\definecolor{slateB}{HTML}{4B5B6E}
\definecolor{slateC}{HTML}{93A6BB}
\definecolor{slateD}{HTML}{2E3B49}
\definecolor{keepblue}{HTML}{E4F0FF}
\definecolor{trustgreen}{HTML}{E6F7E6}

% ---------- page border ----------
\AddToShipoutPictureBG{%
 \begin{tikzpicture}[remember picture,overlay]
  \draw[line width=2.2pt,hudnavy,rounded corners=6pt]
   ([shift={(0.3in,0.3in)}]current page.south west) rectangle ([shift={(-0.3in,-0.3in)}]current page.north east);
  \draw[line width=0.6pt,hudorange,rounded corners=4pt]
   ([shift={(0.37in,0.37in)}]current page.south west) rectangle ([shift={(-0.37in,-0.37in)}]current page.north east);
 \end{tikzpicture}}

% ---------- the guide robot (original design) ----------
\newcommand{\tallybot}[1][hudgreen]{%
 \begin{tikzpicture}[scale=0.42,line join=round,line width=0.4pt,baseline=(current bounding box.center)]
  \coordinate (T) at (0,1.3); \coordinate (UL) at (-1,0.6); \coordinate (UR) at (1,0.6);
  \coordinate (LL) at (-0.9,-0.6); \coordinate (LR) at (0.9,-0.6); \coordinate (B) at (0,-1.1); \coordinate (C) at (0,0.05);
  \fill[hudorange] (-0.28,-1.12)--(0.28,-1.12)--(0,-1.85)--cycle;
  \fill[hudamber] (-0.14,-1.12)--(0.14,-1.12)--(0,-1.5)--cycle;
  \fill[hudorange] (-1,0.5)--(-1.6,0.1)--(-1,-0.25)--cycle;
  \fill[hudorange] (1,0.5)--(1.6,0.1)--(1,-0.25)--cycle;
  \fill[slateC] (T)--(UL)--(C)--cycle;
  \fill[slateA] (T)--(UR)--(C)--cycle;
  \fill[slateB] (UL)--(LL)--(C)--cycle;
  \fill[slateC!70!slateA] (UR)--(LR)--(C)--cycle;
  \fill[slateD] (LL)--(B)--(C)--cycle;
  \fill[slateB!80!black] (LR)--(B)--(C)--cycle;
  \draw[slateD] (T)--(UL)--(LL)--(B)--(LR)--(UR)--cycle;
  \draw[slateD,line width=0.9pt] (T)--(0,1.85);
  \fill[hudorange] (0,1.9) circle (0.14);
  \fill[black] (0,0.1) circle (0.44);
  \fill[#1] (0,0.1) circle (0.29);
  \fill[white] (0.09,0.2) circle (0.09);
 \end{tikzpicture}}

% ---------- HUD-style aside ----------
\newtcolorbox{hudbox}{enhanced,colback=hudbg,colframe=hudbg,boxrule=0pt,arc=0pt,
 left=6pt,right=6pt,top=5pt,bottom=5pt,
 frame code={\path[draw=hudgreen,line width=1pt,fill=hudbg]
  ([xshift=7pt]frame.south west)--([xshift=-7pt]frame.south east)--([yshift=7pt]frame.south east)
  --([yshift=-7pt]frame.north east)--([xshift=-7pt]frame.north east)--([xshift=7pt]frame.north west)
  --([yshift=-7pt]frame.north west)--([yshift=7pt]frame.south west)--cycle;},
 interior code={}}
\newenvironment{tally}[2][hudgreen]{%
 \par\medskip\noindent\begin{hudbox}%
 \begin{minipage}[c]{1.45cm}\centering\tallybot[#1]\end{minipage}\hfill
 \begin{minipage}[c]{\dimexpr\linewidth-2.1cm\relax}\ttfamily\footnotesize\color{#1}\textbf{MULTIX\,//\,#2}\\\ignorespaces}%
 {\end{minipage}\end{hudbox}\par\medskip}

\newcommand{\owes}[1]{\colorbox{hudbg}{\ttfamily\color{hudamber}\footnotesize\ OWES:~#1\ }}
\newcommand{\paidstamp}{\colorbox{hudbg}{\ttfamily\color{hudgreen}\footnotesize\ OWES:~0\ }}
\newcommand{\lightgreen}{\colorbox{hudbg}{\ttfamily\color{hudgreen}\footnotesize\ \(\bullet\) HAS A REASON\ }}
\newcommand{\lightamber}{\colorbox{hudbg}{\ttfamily\color{hudamber}\footnotesize\ \(\bullet\) ONLY SEEN SO FAR\ }}
\newcommand{\lightred}{\colorbox{hudbg}{\ttfamily\color{hudred}\footnotesize\ \(\bullet\) REFUTED\ }}

% ---------- content boxes ----------
\newtcolorbox{storybox}{enhanced,breakable,colback=paper,colframe=inkbrown,boxrule=1.6pt,arc=3mm,
 drop fuzzy shadow=black!25,left=8pt,right=8pt,top=10pt,bottom=7pt,
 attach boxed title to top left={xshift=8mm,yshift*=-\tcboxedtitleheight/2},
 boxed title style={colback=inkbrown,colframe=inkbrown,arc=1.5mm,boxrule=0pt},
 fonttitle=\bfseries\sffamily\footnotesize\color{white},title={STORY}}
\newtcolorbox{tracebox}{enhanced,breakable,colback=hudbg,colframe=hudgreen!80!black,boxrule=1pt,sharp corners,
 left=6pt,right=6pt,top=4pt,bottom=4pt,
 attach boxed title to top left={xshift=6mm,yshift*=-\tcboxedtitleheight/2},
 boxed title style={colback=hudgreen!80!black,colframe=hudgreen!80!black,sharp corners,boxrule=0pt},
 fonttitle=\bfseries\ttfamily\footnotesize\color{hudbg},title={TRACE}}
\newtcolorbox{keepbox}{enhanced,colback=keepblue,colframe=hudnavy,boxrule=1.4pt,arc=3mm,
 left=8pt,right=8pt,top=9pt,bottom=6pt,
 attach boxed title to top left={xshift=8mm,yshift*=-\tcboxedtitleheight/2},
 boxed title style={colback=hudnavy,colframe=hudnavy,arc=1.5mm,boxrule=0pt},
 fonttitle=\bfseries\sffamily\footnotesize\color{white},title={WHAT DID WE KEEP?}}
\newtcolorbox{trustbox}{enhanced,colback=trustgreen,colframe=green!45!black,boxrule=1.4pt,arc=3mm,
 left=8pt,right=8pt,top=9pt,bottom=6pt,
 attach boxed title to top left={xshift=8mm,yshift*=-\tcboxedtitleheight/2},
 boxed title style={colback=green!45!black,colframe=green!45!black,arc=1.5mm,boxrule=0pt},
 fonttitle=\bfseries\sffamily\footnotesize\color{white},title={WHAT MAKES THIS TRUSTWORTHY?}}
\newtcolorbox{braveaside}{enhanced,colback=white,colframe=hudorange,boxrule=1.2pt,arc=3mm,
 borderline={0.4pt}{1.5pt}{hudorange,dashed},left=8pt,right=8pt,top=8pt,bottom=5pt,
 attach boxed title to top right={xshift=-8mm,yshift*=-\tcboxedtitleheight/2},
 boxed title style={colback=hudorange,colframe=hudorange,arc=1.5mm,boxrule=0pt},
 fonttitle=\bfseries\sffamily\footnotesize\color{white},title={FOR BRAVE READERS}}

% margin sticky note
\newcommand{\mnote}[1]{\marginnote{\begin{tcolorbox}[enhanced,colback=yellow!22,colframe=hudamber!70!black,
 boxrule=0.7pt,arc=1.2mm,left=3pt,right=3pt,top=2pt,bottom=2pt,width=\marginparwidth,rotate=-1.5,
 fontupper=\scriptsize\sffamily\raggedright]#1\end{tcolorbox}}}

% chapter banner
\newcommand{\chapterbanner}[3]{%
 \begin{tikzpicture}
  \node[inner sep=0pt] (b) {\begin{minipage}{\textwidth}\centering
   {\sffamily\scriptsize\bfseries\color{hudnavy}#1}\\[2pt]
   {\sffamily\small\color{hudorange}\bfseries CHAPTER #2}\\[3pt]
   {\sffamily\LARGE\bfseries\color{inkbrown}#3}
  \end{minipage}};
  \draw[hudnavy,line width=1.6pt,decorate,decoration={zigzag,segment length=6pt,amplitude=1.2pt}]
   ($(b.south west)+(-0.1,-0.25)$)--($(b.south east)+(0.1,-0.25)$);
  \draw[hudnavy,line width=1.6pt,decorate,decoration={zigzag,segment length=6pt,amplitude=1.2pt}]
   ($(b.north west)+(-0.1,0.3)$)--($(b.north east)+(0.1,0.3)$);
 \end{tikzpicture}\par\medskip}

% puzzle page banner
\newcommand{\puzzlebanner}[2]{%
 \begin{tcolorbox}[enhanced,colback=hudnavy,colframe=hudorange,boxrule=2pt,arc=4mm,
  left=8pt,right=8pt,top=5pt,bottom=5pt,fontupper=\sffamily\color{white}]
  {\scriptsize\bfseries\color{hudamber}PUZZLE PAGE}\\[1pt]{\large\bfseries #1}\\[1pt]{\footnotesize #2}
 \end{tcolorbox}\par\smallskip}


\definecolor{balloonorange}{HTML}{FF9F43}
\definecolor{balloongreen}{HTML}{6CCB5F}
\definecolor{balloonblue}{HTML}{5BA8FF}
\definecolor{balloonpurple}{HTML}{A98BF0}
\definecolor{balloonpink}{HTML}{FF8FC0}

% additions for converted books
\newlength{\stw}
\newcommand{\st}[1]{\settowidth{\stw}{#1}\leavevmode\rlap{\textcolor{red!70!black}{\rule[0.5ex]{\stw}{0.6pt}}}#1}
\usepackage{longtable}
\usepackage{array}
\usepackage{calc}
\providecommand{\tightlist}{\setlength{\itemsep}{0pt}\setlength{\parskip}{0pt}}
\newcommand{\pnum}[1]{\raisebox{0.1ex}{\textcolor{hudorange}{\sffamily\bfseries #1.}}}
\newcommand{\stephead}[2]{\par\medskip\noindent\tikz[baseline=(n.base)]\node[fill=hudnavy,text=white,font=\sffamily\bfseries\scriptsize,inner sep=3pt,rounded corners=2pt](n){STEP #1};\ \ {\sffamily\bfseries\large #2}\par\nopagebreak\smallskip}
\newcommand{\chip}[1]{\par\medskip\noindent\colorbox{hudbg}{\ttfamily\color{hudgreen}\scriptsize\ #1\ }\par\nopagebreak\smallskip}
\newtcolorbox{strategybox}{enhanced,colback=white,colframe=hudnavy,boxrule=1pt,arc=2mm,
 left=8pt,right=8pt,top=6pt,bottom=4pt,
 attach boxed title to top left={xshift=8mm,yshift*=-\tcboxedtitleheight/2},
 boxed title style={colback=hudnavy,colframe=hudnavy,arc=1.5mm,boxrule=0pt},
 fonttitle=\bfseries\sffamily\scriptsize\color{white},title={SPLIT $\cdot$ BALANCE $\cdot$ CORRECT}}
\newtcolorbox{predictbox}{enhanced,colback=yellow!12,colframe=hudamber!80!black,boxrule=1pt,arc=2mm,
 left=8pt,right=8pt,top=7pt,bottom=4pt,
 attach boxed title to top left={xshift=8mm,yshift*=-\tcboxedtitleheight/2},
 boxed title style={colback=hudamber!80!black,colframe=hudamber!80!black,arc=1.5mm,boxrule=0pt},
 fonttitle=\bfseries\sffamily\scriptsize\color{white},title={PREDICT FIRST}}
\newtcolorbox{practicebox}{enhanced,breakable,colback=white,colframe=hudorange,boxrule=1.6pt,arc=3mm,
 left=8pt,right=8pt,top=9pt,bottom=6pt,
 before upper={\setlist[enumerate,1]{label=\pnum{\arabic*},leftmargin=2.4em,itemsep=0.45em,topsep=0pt}},
 attach boxed title to top left={xshift=8mm,yshift*=-\tcboxedtitleheight/2},
 boxed title style={colback=hudorange,colframe=hudorange,arc=1.5mm,boxrule=0pt},
 fonttitle=\bfseries\sffamily\footnotesize\color{white},title={PRACTICE}}
\renewtcolorbox{trustbox}[1][WHAT MAKES THIS TRUSTWORTHY?]{enhanced,colback=trustgreen,colframe=green!45!black,boxrule=1.4pt,arc=3mm,
 left=8pt,right=8pt,top=9pt,bottom=6pt,
 attach boxed title to top left={xshift=8mm,yshift*=-\tcboxedtitleheight/2},
 boxed title style={colback=green!45!black,colframe=green!45!black,arc=1.5mm,boxrule=0pt},
 fonttitle=\bfseries\sffamily\scriptsize\color{white},title={#1}}
\newtcolorbox{bridgebox}{enhanced,colback=hudnavy,colframe=hudnavy,boxrule=0pt,arc=2mm,
 left=8pt,right=8pt,top=6pt,bottom=5pt,fontupper=\sffamily\scriptsize\color{white},
 before upper={\textcolor{hudamber}{\bfseries BRIDGE\ \ }}}
\newcommand{\bigmultix}{\scalebox{3.2}{\tallybot}}

\newtcolorbox{unitbox}[1]{enhanced,colback=white,colframe=hudnavy,boxrule=1.2pt,arc=2.5mm,
 left=8pt,right=6pt,top=9pt,bottom=5pt,
 attach boxed title to top left={xshift=8mm,yshift*=-\tcboxedtitleheight/2},
 boxed title style={colback=hudnavy,colframe=hudnavy,arc=1.5mm,boxrule=0pt},
 fonttitle=\bfseries\sffamily\scriptsize\color{white},title={#1}}
\begin{document}
\thispagestyle{empty}
\vspace*{1.1cm}
\begin{center}
{\sffamily\bfseries\footnotesize\color{hudorange} CURRICULUM OVERVIEW}\\[6pt]
{\sffamily\bfseries\fontsize{26}{28}\selectfont\color{inkbrown} One Year,\\ Five Series}\\[10pt]
{\sffamily\small\color{hudnavy} Forty-five short books for ages 9 to 11}\\[1.1cm]
\bigmultix\\[1.1cm]
\begin{tally}[hudgreen]{HELLO}
Five subjects, nine books each, read side by side in nine units. I ride along in every one.
\end{tally}
\end{center}
\newpage

\chapterbanner{OVERVIEW}{1}{The Five Series}

\begin{storybox}
Each series has a central question its nine books return to. Every book teaches the same habit: say what you actually have, and what would change your mind.
\end{storybox}

\vspace{-4pt}\stephead{1}{Mathematics}
\textbf{What stays true when we transform a problem?} Nine books built around one running problem, $48\times52$, and the circle, pop, redraw and refusal notation. Recurring questions: \textit{What did we keep?} and \textit{What makes this trustworthy?}

\stephead{2}{Science}
\textbf{How We Find Things Out.} What evidence supports an explanation? Life, physical and Earth science, with household experiments. Recurring questions: \textit{What did we actually observe?} and \textit{What would change our minds?} A third, \textit{What remains unknown?}, appears only when something truly is.

\stephead{3}{Spelling and Grammar}
\textbf{What Words Can Do.} What happens to meaning when we change an expression? Spelling, grammar, meaning and revision. Recurring questions: \textit{What changed when we changed the words?} and \textit{How can we tell which version works?} Standard written English is taught as a convention for certain situations, never as a correction of home dialects.

\stephead{4}{Ecology}
\textbf{The World Between Living Things.} How do relationships let living systems persist and change? Five places, a pond, woodland, meadow, shoreline and garden, recur across the books, with Atlantic Canadian examples. Recurring questions: \textit{Who or what depends on this?}, \textit{What changed, and what remained?} and \textit{What haven't we counted?}

\stephead{5}{Problem Solving}
\textbf{A history of problems, from prehistory to the modern world.} What did people try, under which constraints, and what did the solution change? Recurring questions: \textit{What was the problem?}, \textit{What could they try?} and \textit{What did the solution change?}

\newpage
\chapterbanner{OVERVIEW}{2}{The Year Map}

\begin{tally}[hudgreen]{PACING}
Nine units of four weeks each makes 36 weeks. In every unit, book $n$ of each series is read together. That is about 40 short chapters per unit, or 10 a week, two per series.
\end{tally}

\begin{unitbox}{UNIT 1 \quad Weeks 1 to 4}
\begin{tabular}{@{}p{1.3in}p{2.5in}@{}}
\textcolor{hudorange}{\textbf{Math}} & \textit{A Useful Wrong Answer}\\
\textcolor{hudorange}{\textbf{Science}} & \textit{The Thing We Didn't See}\\
\textcolor{hudorange}{\textbf{Language}} & \textit{The Letter That Changed Everything}\\
\textcolor{hudorange}{\textbf{Ecology}} & \textit{The Forest Under Our Feet}\\
\textcolor{hudorange}{\textbf{Problems}} & \textit{The Stone That Became a Tool}\\
\end{tabular}
\end{unitbox}

\begin{unitbox}{UNIT 2 \quad Weeks 5 to 8}
\begin{tabular}{@{}p{1.3in}p{2.5in}@{}}
\textcolor{hudorange}{\textbf{Math}} & \textit{The Work That Vanishes}\\
\textcolor{hudorange}{\textbf{Science}} & \textit{The Experiment That Changed Two Things}\\
\textcolor{hudorange}{\textbf{Language}} & \textit{The Word Built from Pieces}\\
\textcolor{hudorange}{\textbf{Ecology}} & \textit{The Pond That Changed Overnight}\\
\textcolor{hudorange}{\textbf{Problems}} & \textit{Keeping the Fire}\\
\end{tabular}
\end{unitbox}

\begin{unitbox}{UNIT 3 \quad Weeks 9 to 12}
\begin{tabular}{@{}p{1.3in}p{2.5in}@{}}
\textcolor{hudorange}{\textbf{Math}} & \textit{Two Boxes, One Answer}\\
\textcolor{hudorange}{\textbf{Science}} & \textit{The Measurement That Moved}\\
\textcolor{hudorange}{\textbf{Language}} & \textit{The Sentence That Turned Around}\\
\textcolor{hudorange}{\textbf{Ecology}} & \textit{The Flower That Needed a Visitor}\\
\textcolor{hudorange}{\textbf{Problems}} & \textit{The Food That Had to Last}\\
\end{tabular}
\end{unitbox}

\begin{unitbox}{UNIT 4 \quad Weeks 13 to 16}
\begin{tabular}{@{}p{1.3in}p{2.5in}@{}}
\textcolor{hudorange}{\textbf{Math}} & \textit{The Order of Discovery}\\
\textcolor{hudorange}{\textbf{Science}} & \textit{The Model That Almost Worked}\\
\textcolor{hudorange}{\textbf{Language}} & \textit{The Comma That Changed the Story}\\
\textcolor{hudorange}{\textbf{Ecology}} & \textit{The Journey Without a Map}\\
\textcolor{hudorange}{\textbf{Problems}} & \textit{The River That Wouldn't Obey}\\
\end{tabular}
\end{unitbox}

\begin{unitbox}{UNIT 5 \quad Weeks 17 to 20}
\begin{tabular}{@{}p{1.3in}p{2.5in}@{}}
\textcolor{hudorange}{\textbf{Math}} & \textit{Reversing the Calculation}\\
\textcolor{hudorange}{\textbf{Science}} & \textit{The History Inside the Seed}\\
\textcolor{hudorange}{\textbf{Language}} & \textit{Two Meanings, One Sentence}\\
\textcolor{hudorange}{\textbf{Ecology}} & \textit{The Meadow That Became a Forest}\\
\textcolor{hudorange}{\textbf{Problems}} & \textit{Finding an Island Without a Map}\\
\end{tabular}
\end{unitbox}

\begin{unitbox}{UNIT 6 \quad Weeks 21 to 24}
\begin{tabular}{@{}p{1.3in}p{2.5in}@{}}
\textcolor{hudorange}{\textbf{Math}} & \textit{What Must Remain Unresolved}\\
\textcolor{hudorange}{\textbf{Science}} & \textit{The Pond That Changed Its Mind}\\
\textcolor{hudorange}{\textbf{Language}} & \textit{The Rewrite That Lost Something}\\
\textcolor{hudorange}{\textbf{Ecology}} & \textit{The Shoreline That Wouldn't Stay Still}\\
\textcolor{hudorange}{\textbf{Problems}} & \textit{The Water That Had to Climb}\\
\end{tabular}
\end{unitbox}

\begin{unitbox}{UNIT 7 \quad Weeks 25 to 28}
\begin{tabular}{@{}p{1.3in}p{2.5in}@{}}
\textcolor{hudorange}{\textbf{Math}} & \textit{The Choice Before the Calculation}\\
\textcolor{hudorange}{\textbf{Science}} & \textit{The Invisible Part of the Picture}\\
\textcolor{hudorange}{\textbf{Language}} & \textit{The Perfectly Wrong Sentence}\\
\textcolor{hudorange}{\textbf{Ecology}} & \textit{The Forest That Grew Back Differently}\\
\textcolor{hudorange}{\textbf{Problems}} & \textit{The Book That Took Too Long}\\
\end{tabular}
\end{unitbox}

\begin{unitbox}{UNIT 8 \quad Weeks 29 to 32}
\begin{tabular}{@{}p{1.3in}p{2.5in}@{}}
\textcolor{hudorange}{\textbf{Math}} & \textit{The Boundary of a Step}\\
\textcolor{hudorange}{\textbf{Science}} & \textit{Two Explanations, One Result}\\
\textcolor{hudorange}{\textbf{Language}} & \textit{The Spelling That Remembered}\\
\textcolor{hudorange}{\textbf{Ecology}} & \textit{The Species Nobody Counted}\\
\textcolor{hudorange}{\textbf{Problems}} & \textit{The Map That Found a Pattern}\\
\end{tabular}
\end{unitbox}

\begin{unitbox}{UNIT 9 \quad Weeks 33 to 36}
\begin{tabular}{@{}p{1.3in}p{2.5in}@{}}
\textcolor{hudorange}{\textbf{Math}} & \textit{Solving by Changing the Question}\\
\textcolor{hudorange}{\textbf{Science}} & \textit{The Answer We Haven't Earned}\\
\textcolor{hudorange}{\textbf{Language}} & \textit{The Message That Finally Worked}\\
\textcolor{hudorange}{\textbf{Ecology}} & \textit{The Garden That Could Not Be Finished}\\
\textcolor{hudorange}{\textbf{Problems}} & \textit{The Solution Nobody Planned}\\
\end{tabular}
\end{unitbox}

\newpage
\chapterbanner{OVERVIEW}{3}{A Week With Multix}

\begin{strategybox}
Two chapters from each series every week. Each chapter is one sitting of about 20 minutes, plus its puzzle page when it has one.
\end{strategybox}

\begin{center}\sffamily\small
\begin{tabular}{@{}lll@{}}
\toprule
\textbf{Day} & \textbf{First sitting} & \textbf{Second sitting}\\
\midrule
Monday & Mathematics & Science\\
Tuesday & Spelling and Grammar & Ecology\\
Wednesday & Problem Solving & Mathematics\\
Thursday & Science & Spelling and Grammar\\
Friday & Ecology & Problem Solving\\
\bottomrule
\end{tabular}
\end{center}

\stephead{1}{Shared habits}
Three habits show up under different names in all five series.
\begin{itemize}
\item \textbf{Label the claim.} Mathematics sorts claims into \textit{has a reason}, \textit{only seen so far} and \textit{refuted}. Science separates observed from inferred. Problem Solving separates evidence, reconstruction and imagined scenario.
\item \textbf{Change one thing.} A fair experiment, a rewritten sentence, a rearranged sum and a changed pond all ask what moved and what stayed.
\item \textbf{Say what is missing.} An owed correction, an unresolved question, an uncounted species and an unrecorded intention are all things the notebook should name.
\end{itemize}

\stephead{2}{Puzzle pages and answer keys}
Every chapter ends with a puzzle page: spot the slip, color by answer, balloon message, and others added as each series needs them. Every number, measurement, graph and key is generated from checked source data and verified independently before it is typeset. Illustrations never carry the math.

\newpage
\chapterbanner{OVERVIEW}{4}{Open Points}

\stephead{1}{Overlap to settle}
The Science book \textit{The Pond That Changed Its Mind} (Unit 6) and the Ecology book \textit{The Pond That Changed Overnight} (Unit 2) cover close to the same ground, and their titles are nearly the same. Options: keep both and give the Science book a measurement-and-feedback angle, retitle one, or move the Science pond to a different unit from the Ecology pond.

\stephead{2}{Shape of the titles}
Of the 45 titles, 23 follow the pattern \textit{The ... That ...}. The pattern is pleasant in a small dose and tiring across a shelf. A few of the Science, Language and Ecology titles could be reworked without losing anything.

\stephead{3}{Counts to confirm}
Mathematics has 73 chapters, eight per book and nine in Book 3. The other four series are planned at about eight chapters a book, with Problem Solving's six to eight depending on how many episodes each theme holds. The pacing above assumes eight.

\stephead{4}{Facts before stories}
Problem Solving is the one series where invented intentions could pass as established fact. Each historical chapter is drafted only after its case study is checked against primary evidence and reliable scholarship, and the three kinds of statement, evidence, reconstruction and imagined scenario, are marked on the page.

\stephead{5}{Order of work}
Mathematics is already drafted and one book is typeset. A sensible order is to typeset the rest of Mathematics, then draft and typeset each other series one unit at a time so every series stays in step with the year map.
\end{document}
